\section{Mathematical construction on curvilinear elements}
\label{sec:curvilinear-mathematics}
The affine implementation relies on two simplifications: the Jacobian is
constant within an element, and differentiation stays inside the represented
polynomial space. A curved map removes both simplifications. Our extension
keeps the physical projection idea of OFDG, but assembles its operators for
each curved element and projects each derivative back into the solution space.
This section derives that construction; Appendix~\ref{sec:ofdg-implementation}
describes its coefficient storage and evaluation.

\subsection{Why the affine construction is insufficient}
Let $\boldsymbol x=T_K(\boldsymbol\xi)$ be a smooth, invertible,
orientation-preserving element map, and write
\begin{equation}
 J_K(\boldsymbol\xi)=\frac{\partial T_K}{\partial\boldsymbol\xi},
 \qquad w_K(\boldsymbol\xi)=\det J_K(\boldsymbol\xi)>0.
 \label{eq:curved-jacobian}
\end{equation}
For $v=\widehat v\circ T_K^{-1}$, change of variables and the chain rule give
\begin{align}
 \int_K v\,dx&=\int_{\widehat K}\widehat v\,w_K\,d\xi,
 \label{eq:curved-change-variables}\\
 (\nabla_xv)\circ T_K&=J_K^{-\mathsf T}\nabla_\xi\widehat v.
 \label{eq:curved-chain-rule}
\end{align}
The physical inner product therefore has a nonconstant weight. Even if
$\widehat v$ and $T_K$ are polynomial, the inverse Jacobian generally makes
$\partial_{x_j}v$ fall outside the mapped finite element space. Evaluating a
single inverse Jacobian at the element centre cannot represent this variation.

A one-dimensional example makes the obstruction explicit. On $[0,1]$, take
\begin{equation}
 x=T(\xi)=\xi+\varepsilon\xi(1-\xi),\quad
 0<|\varepsilon|<1,\qquad \widehat v(\xi)=\xi.
\end{equation}
Although the reference state is linear,
\begin{equation}
 (\partial_xv)\circ T=\frac{1}{1+\varepsilon-2\varepsilon\xi},\qquad
 (\partial_x^2v)\circ T=
 \frac{2\varepsilon}{(1+\varepsilon-2\varepsilon\xi)^3}.
 \label{eq:curved-derivative-example}
\end{equation}
Neither expression is generally a finite-degree reference polynomial. The
issue is already present for an ordinary polynomial geometry map; it does not
require NURBS.

\subsection{Physical mass matrices and nested projections}
Use the mapped spaces $V^r(K)$ from \cref{eq:ofdg-space}, all with the same
map $T_K$. Let $\widehat\phi_i$ span the order-$k$ reference space and
$\widehat\psi_a^{(r)}$ its order-$r$ subspace. Their physical mass and coupling
matrices are
\begin{align}
 (M_k^K)_{ij}&=\int_{\widehat K}\widehat\phi_i\widehat\phi_j w_K\,d\xi,\\
 (M_r^K)_{ab}&=\int_{\widehat K}
       \widehat\psi_a^{(r)}\widehat\psi_b^{(r)}w_K\,d\xi,\\
 (B_r^K)_{ai}&=\int_{\widehat K}
       \widehat\psi_a^{(r)}\widehat\phi_iw_K\,d\xi.
 \label{eq:curved-weighted-masses}
\end{align}
If $U$ represents $u_h$, its physical $L^2$ projection into $V^r(K)$ has
lower-order coefficients $C_r$ determined by
\begin{equation}
 M_r^K C_r=B_r^K U.
 \label{eq:curved-projection-solve}
\end{equation}
This is the usual finite element projection equation, now with the physical
Jacobian weight retained. To express the projected function in the original
order-$k$ basis, solve $M_k^K U_r=(B_r^K)^{\mathsf T}C_r$. Thus the matrix
representing $\Pi_K^r$ in that basis is
\begin{equation}
 P_r^K=(M_k^K)^{-1}(B_r^K)^{\mathsf T}(M_r^K)^{-1}B_r^K.
 \label{eq:curved-projector-matrix}
\end{equation}
The inverse notation denotes the associated linear solves. Nesting makes
this lifting exact as a function-space operation. With exact integrals,
\begin{equation}
 P_r^KP_s^K=P_{\min(r,s)}^K,\qquad
 (P_r^K)^{\mathsf T}M_k^K=M_k^KP_r^K.
 \label{eq:curved-projector-identities}
\end{equation}
Consequently the physical orthogonal shells and the frozen exponential
attenuation in \cref{eq:ofdg-exact-decay} remain available. No orthogonal
storage basis is required.

Mean preservation follows directly from testing the projection equation with
one: $\int_K(\Pi_K^r u_h-u_h)\,dx=0$. In coefficient form, let
$b_i=\int_K\phi_i\,dx$; then $b^{\mathsf T}P_r^K=b^{\mathsf T}$. Every
nonconstant shell has zero $b$-weighted integral, so damping it leaves
$b^{\mathsf T}U$ unchanged component by component. On affine elements,
$w_K$ is constant and cancels from \cref{eq:curved-projector-matrix}, recovering
the reusable reference projector. On curved elements it cannot be cancelled.

\subsection{Projecting a physical derivative}
We seek $z_h=\mathscr D_{K,j}u_h\in V^k(K)$ satisfying
\begin{equation}
 (z_h,v_h)_K=(\partial_{x_j}u_h,v_h)_K
 \quad\text{for every }v_h\in V^k(K).
 \label{eq:curved-derivative-weak}
\end{equation}
This is an $L^2$ projection, not a new PDE solve or a finite-difference
approximation. Inserting the basis expansions and
\cref{eq:curved-chain-rule} gives
\begin{align}
 M_k^K Z&=G_{K,j}U,\qquad D_{K,j}=(M_k^K)^{-1}G_{K,j},\\
 (G_{K,j})_{ab}&=\int_{\widehat K}\widehat\phi_a
  \left[\sum_{s=1}^d(J_K^{-1})_{sj}
                   \partial_{\xi_s}\widehat\phi_b\right]w_K\,d\xi.
 \label{eq:curved-derivative-assembly}
\end{align}
The geometry is evaluated inside the integral. Only first reference
basis derivatives and first derivatives of the map are needed. A constant
has zero right-hand side, hence zero projected derivative. For affine maps,
physical differentiation stays in the supported space and this projection
recovers the existing exact polynomial derivative matrix.

In the one-dimensional example \cref{eq:curved-derivative-example}, the
projected first derivative $z_h$ is characterized by
\begin{equation}
 \int_0^1\widehat z_h\widehat\phi_a
       (1+\varepsilon-2\varepsilon\xi)\,d\xi
 =\int_0^1\widehat\phi_a\,d\xi.
\end{equation}
The Jacobian cancels on the right, but remains in the mass matrix on the left.
This illustrates why a physical projection is different from an unweighted
reference projection.

\subsection{Recursive higher derivatives: the precise approximation}
The derivative graph reuses these first-derivative matrices:
\begin{equation}
 U^{(0)}=U,\qquad U^{(\alpha)}=D_{K,j}U^{(\alpha-e_j)},
 \label{eq:curved-recursion}
\end{equation}
where $e_j$ selects the coordinate direction assigned to that graph edge.
For instance, the second operation in directions $i$ then $j$ represents
$\Pi_K^k\partial_{x_j}(\Pi_K^k\partial_{x_i}u_h)$. It generally differs
from projecting the exact second derivative once. Writing $\Pi=\Pi_K^k$,
\begin{equation}
 \Pi\partial_{x_j}(\Pi\partial_{x_i}u_h)
 -\Pi\partial_{x_j}\partial_{x_i}u_h
 =\Pi\partial_{x_j}(\Pi-I)\partial_{x_i}u_h.
 \label{eq:curved-recursion-defect}
\end{equation}
The right-hand side is the differentiated error from the previous projection.
It explains both the approximation and why projected mixed derivatives need
not commute. The implementation fixes an order of application; it does not
average all orders or claim exact higher chain-rule derivatives. Second- and
third-derivative convergence must therefore be assessed using
manufactured solutions. Projection alone is not a proof of their convergence.

\subsection{Curved faces, normal directions, and length scales}
For a reference face with outward unit normal $\widehat{\boldsymbol n}$, the
physical surface factor and outward normal are
\begin{equation}
 s_F=\|w_KJ_K^{-\mathsf T}\widehat{\boldsymbol n}\|,
 \qquad \boldsymbol n_K=
 \frac{w_KJ_K^{-\mathsf T}\widehat{\boldsymbol n}}{s_F},
 \qquad ds=s_F\,d\widehat s.
 \label{eq:curved-face-metric}
\end{equation}
These quantities vary along a curved face. For face quadrature weights
$\omega_q$, use
\begin{equation}
 |F|_Q=\sum_q\omega_qs_{F,q},\qquad
 \widetilde\omega_q=\frac{\omega_qs_{F,q}}{|F|_Q},\qquad
 \sum_q\widetilde\omega_q=1.
 \label{eq:curved-face-normalization}
\end{equation}
Hence the implemented derivative-jump average is
\begin{equation}
 J_{F,c}^{(l)}\approx\sum_{|\alpha|=l}\sum_q\widetilde\omega_q
 \left| (\mathscr D_{K^-}^{\alpha}u_{h,c}^-)(x_q)
       -(\mathscr D_{K^+}^{\alpha}u_{h,c}^+)(x_q)\right|.
\end{equation}
Both traces are evaluated at the same physical point, using their respective
element maps. A constant jump has the same normalized magnitude on a straight
or curved face. The one-sided wave radii use $\boldsymbol n_K(x_q)$ at these
points, rather than a single centre normal.

Volume quadrature gives $|K|_Q=\sum_q\omega_qw_K(\xi_q)$ and the OFDG length
$h_K=(|K|_Q/|\widehat K|)^{1/d}$. The global mean and amplitude calculation
uses the same physical volume interpretation. KXRCF uses its separately defined
sampled radius and restricts its signed face integral to quadrature points
with incoming transport. Its denominator uses the corresponding physical
inflow measure. Thus sensing, normalization, and propagation directions all
refer to the curved physical geometry.

The algebra above uses exact integrals. The implementation assembles them
with positive, overintegrated quadrature. Consistent quadrature on the nested
spaces gives the analogous discrete projection identities, up to linear-solve
roundoff; it preserves the quadrature-defined physical means. Agreement with
exact integrals, derivative accuracy, and geometric sampling invariance still
require quadrature and refinement checks. This solution-space
projection does not convert rational NURBS geometry to polynomial geometry;
that separate, explicit geometry approximation retains the support boundary
stated in \cref{sec:support-boundary}.
